% ============ §20 UFDs (原书 pp.161-169) ============
\section{UFDs}\label{sec:20}\index{UFD, section 20}\index{factorial ring, section 20}

This section treats UFDs, which we have already touched on in
\S\ref{sec:1}; note that the Bourbaki terminology for UFD is `factorial ring'. First of
all, we have the following criterion for Noetherian rings.

\begin{theorem}\label{thm:20.1}
A Noetherian integral domain $A$ is a UFD if and only if
every height 1 prime ideal is principal.
\end{theorem}

\begin{proof}
Proof of `only if'. Suppose that $A$ is a UFD and that $P$ is a height 1
prime ideal. Take any non-zero $a \in P$, and express $a$ as a product of prime
elements, $a = \prod \pi_i$. Then at least one of the $\pi_i$ belongs to $P$; if $\pi_i \in P$ then
$(\pi_i) \subset P$, but $(\pi_i)$ is a non-zero prime ideal and $\htop P = 1$, hence $P = (\pi_i)$.

Proof of `if'. Since $A$ is Noetherian, every element $a \in A$ which is neither
0 nor a unit can be written as a product of finitely many irreducibles.
Hence it will be enough to prove that an irreducible element $a$ is a prime
element. Let $P$ be a minimal prime divisor of $(a)$; then by the principal
ideal theorem (Theorem~\ref{thm:13.9}), $\htop P = 1$, so that by assumption we can
write $P = (b)$. Thus $a = bc$, and since $a$ is irreducible, $c$ is a unit, so that
$(a) = (b) = P$, and $a$ is a prime element.
\end{proof}

\begin{theorem}\label{thm:20.2}
Let $A$ be a Noetherian integral domain, $\Gamma$ a set of prime
elements of $A$, and let $S$ be the multiplicative set generated by $\Gamma$. If $A_S$
is a UFD then so is $A$.
\end{theorem}

\begin{proof}
Let $P$ be a height 1 prime ideal of $A$. If $P \cap S \neq 0$ then $P$ contains
an element $\pi \in \Gamma$, and since $\pi A$ is a non-zero prime ideal we have $P = \pi A$.
If $P \cap S = 0$ then $PA_S$ is a height 1 prime ideal of $A_S$, so that $PA_S = aA_S$
for some $a \in P$. Among all such $a$ choose one such that $aA$ is maximal;
then $a$ is not divisible by any $\pi \in \Gamma$. Now if $x \in P$ we have $sx = ay$ for
some $s \in S$ and $y \in A$. Let $s = \pi_1 \cdots \pi_r$ with $\pi_i \in \Gamma$; then $a \notin \pi_iA$, so
that $y \in \pi_iA$, and an induction on $r$ shows that $y \in sA$, so that $x \in aA$. Hence
$P = aA$.
\end{proof}

\begin{lemma}\label{lem:20.1}
Let $A$ be an integral domain, and $a$ an ideal of $A$ such that
$a \oplus A^n \cong A^{n+1}$; then $a$ is principal.
\end{lemma}

\begin{proof}
Fix the basis $e_0, \ldots, e_n$ of $A^{n+1}$, and viewing $a \oplus A^n \subset A \oplus A^n$, fix
$f_0, \ldots, f_n$ such that $f_0$ is a basis of $A$ and $f_1, \ldots, f_n$ a basis of $A^n$. Then
the isomorphism $\varphi\colon A^{n+1} \to a \oplus A^n$ can be given in the form $\varphi(e_i) =
\sum_{j=0}^{n} a_{ij}f_j$. Write $d_i$ for the $(i, 0)$th cofactor of the matrix $(a_{ij})$, and $d$
for the determinant, so that, since $\varphi$ is injective, $d \neq 0$, and $\sum_i a_{i0}d_i = d$,
$\sum_i a_{ij}d_i = 0$ if $j \neq 0$. Hence if we set $e_0' = \sum_i d_ie_i$, we have $\varphi(e_0') = df_0$.
Moreover, since the image of $\varphi$ includes $f_1, \ldots, f_n$, there exist $e_1', \ldots,
e_n' \in A^{n+1}$ such that $\varphi(e_j') = f_j$. Now define a matrix $(c_{jk})$ by $e_j' = \sum_{k=0}^{n} c_{jk}e_k$
for $j = 0, \ldots, n$ (so $c_{0k} = d_k$). Then we have
\[
  (c_{jk})(a_{ij}) = \begin{pmatrix}
    d & 0 & \cdots & 0 \\
    0 & 1 & \cdots & 0 \\
    \vdots & \vdots & \ddots & \vdots \\
    0 & 0 & \cdots & 1
  \end{pmatrix},
\]
so that by comparing the determinants of both sides we get $\det(c_{jk}) = 1$.
Therefore $e_0', \ldots, e_n'$ is another basis of $A^{n+1}$, and $\varphi(Ae_0') = dAf_0$,
so that $a = dA$.
\end{proof}

Let $K$ be the field of fractions of the integral domain $A$; for a finite
$A$-module $M$, the dimension of $M \otimes_A K$ as a vector space over $K$ is called
the \emph{rank}\index{rank of a module} of $M$. A torsion-free finite $A$-module of rank 1 is isomorphic to
an ideal of $A$. Lemma~\hyperref[lem:20.1]{1} can be formulated as saying that for an integral
domain $A$, a stably free rank 1 module is free (see Ex.~\ref{ex:19.2-ex}). The elementary
proof given above is taken from a lecture by M. Narita in 1971.

\begin{theorem}[Auslander and Buchsbaum~\cite{AB3}]\label{thm:20.3}
A regular local ring is a
UFD.
\end{theorem}

\begin{proof}
Let $(A, m)$ be a regular local ring: the proof works by induction
on $\dim A$. If $\dim A = 0$ then $A$ is a field and therefore (trivially) a UFD.
If $\dim A = 1$ then $A$ is a DVR, and therefore a UFD. We suppose that
$\dim A > 1$ and choose $x \in m - m^2$; then since $xA$ is a prime ideal, applying
Theorem~\hyperref[thm:20.2]{2} to $\Gamma = \{x\}$, we need only show that $A_x$ is a UFD (where
$A_x = A[x^{-1}]$ is as on p.~22). Let $P$ be a height 1 prime ideal of $A_x$, and
set $p = P \cap A$; we have $P = pA_x$. Since $A$ is a regular local ring, the
$A$-module $p$ has an FFR, so that the $A_x$-module $P$ has an FFR. For any
prime ideal $Q$ of $A_x$, the ring $(A_x)_Q = A_{Q \cap A}$ is a regular local ring of
dimension less than that of $A$, so by induction is a UFD. Thus $P_Q$ is free
as an $(A_x)_Q$-module, so that by Theorem~\ref{thm:7.12}, the $A_x$-module $P$ is projective;
hence by Ex.~\ref{ex:19.2-ex}, $P$ is stably free, and therefore by the previous lemma,
$P$ is a principal ideal of $A_x$.
\end{proof}

The above proof is due to Kaplansky. Instead of our Lemma~\hyperref[lem:20.1]{1}, he used
the following more general proposition, which he had previously proved:
if $A$ is an integral domain, and $I_i, J_i$ are ideals of $A$ for $1 \le i \le r$ such that
$\bigoplus_{i=1}^{r} I_i \cong \bigoplus_{i=1}^{r} J_i$, then $I_1 \cdots I_r \cong J_1 \cdots J_r$. This is an interesting property
of ideals, and we have given a proof in Appendix C.

\begin{theorem}\label{thm:20.4}
Let $A$ be a Noetherian integral domain. Then if any finite
$A$-module has an FFR, $A$ is a UFD.
\end{theorem}

\begin{proof}
By Ex.~\ref{ex:19.4-ex}, $A$ is a regular ring. Let $P$ be a height 1 prime ideal
of $A$. Then $A_m$ is a regular local ring for any $m \in \Spec A$, so by the previous
theorem, the ideal $P_m$ is principal, and is therefore a free $A_m$-module. Hence
by Theorem~\ref{thm:7.12}, $P$ is projective. Therefore by Ex.~\ref{ex:19.2-ex}, $P$ is stably free,
and so by Lemma~\hyperref[lem:20.1]{1} is principal.
\end{proof}

Let $A$ be an integral domain; for any two non-zero elements $a, b \in A$,
the notion of greatest common divisor (g.c.d.) and least common multiple
(1.c.m.) are defined as in the ring of integers. That is, $d$ is a g.c.d. of $a$ and
$b$ if $d$ divides both $a$ and $b$, and any element $x$ dividing both $a$ and $b$
divides $d$; and $e$ is an 1.c.m. of $a$ and $b$ if $e$ is divisible by both $a$ and $b$,
and any $y$ divisible by $a$ and $b$ is divisible by $e$; this condition is equivalent
to $(e) = (a) \cap (b)$.

\begin{lemma}\label{lem:20.2}
If an 1.c.m. of $a$ and $b$ exists then so does a g.c.d.
\end{lemma}

\begin{proof}
If $(a) \cap (b) = (e)$ then there exists $d$ such that $ab = ed$. From $e \in (a)$
we get $b \in (d)$ and similarly $a \in (d)$, so that $(a, b) \subset (d)$. Now if $x$ is a common
divisor of $a$ and $b$ then $a = xt$ and $b = xs$, so that $xst$ is a common multiple
of $a, b$, and is hence divisible by $e$. Then from $ed = ab = x \cdot xst$ we get that
$d$ is divisible by $x$. Therefore, $d$ is a g.c.d. of $a$ and $b$.
\end{proof}

\begin{remark}\label{rem:20.1}
If $A$ is a Noetherian integral domain which is not a UFD then
$A$ has an irreducible element $a$ which is not prime. If $xy \in (a)$ but $x \notin (a)$,
$y \notin (a)$ then the only common divisors of $a$ and $x$ are units, so that 1 is a
g.c.d. of $a$ and $x$. However, $xy \in (a) \cap (x)$, but $xy \notin (ax)$, so that $(a) \cap
(x) \neq (ax)$, and there does not exist any 1.c.m. of $a$ and $x$. Thus the converse
of Lemma~\hyperref[lem:20.2]{2} does not hold in general.
\end{remark}

\begin{remark}\label{rem:20.2}
If $A$ is a UFD then an intersection of an arbitrary collection
of principal ideals is again principal (we include (0)). Indeed, if $\bigcap_i a_iA \neq 0$,
then factorise each $a_i$ as a product of primes:
\[
  a_i = u_i \prod_p p^{n_{ip}},
\]
with $u_i$ units, and $p, \pi$ prime elements such that $pA \neq \pi A$ for $p \neq \pi$.
(We could even allow the $a_i$ to be elements of the field of fractions of $A$.)
Then $\bigcap_i a_iA = dA$, where $d = \prod_p p^{\min_i n_{ip}}$.
\end{remark}

\begin{theorem}\label{thm:20.5}
An integral domain $A$ is a UFD if and only if the ascending
chain condition holds for principal ideals, and any two elements of $A$ have
an 1.c.m.
\end{theorem}

\begin{proof}
The `only if' is already known, and we prove the `if'. From the first
condition it follows that every element which is neither 0 nor a unit can
be written as a product of a finite number of irreducible elements, so that
we need only prove that an irreducible element is prime. Let $a$ be an
irreducible element, and let $xy \in (a)$ and $x \notin (a)$. By assumption we can
write $(a) \cap (x) = (z)$; now 1 is a g.c.d. of $a$ and $x$, so that one sees from the
proof of Lemma~\hyperref[lem:20.2]{2} that $(z) = (ax)$, and then $xy \in (a) \cap (x) = (ax)$ implies
that $y \in (a)$. Therefore $(a)$ is prime.
\end{proof}

\begin{theorem}\label{thm:20.6}
Let $A$ be a regular ring and $u, v \in A$. Then $uA \cap vA$ is a
projective ideal.
\end{theorem}

\begin{proof}
$A_m$ is a UFD for every maximal ideal $m$, so that $(uA \cap vA)A_m =
uA_m \cap vA_m$ is a principal ideal, and hence a free module.
\end{proof}

\begin{theorem}\label{thm:20.7}
If $A$ is a UFD then a projective ideal is principal.
\end{theorem}

\begin{proof}
By Theorem~\ref{thm:11.3}, it is equivalent to say that a non-zero ideal $a$ is
projective or invertible. Hence if we set $K$ for the field of fractions of $A$,
then there exist $u_i \in K$ such that $u_ia \subset A$ and $a_i \in a$ such that $\sum u_ia_i = 1$.
We have $a \subset \bigcap u_i^{-1}A$, and conversely if $x \in \bigcap u_i^{-1}A$ then $x =
\sum (xu_i)a_i \in a$, and hence $a = \bigcap u_i^{-1}A$; now since $A$ is a UFD, the intersection
of principal fractional ideals is again principal.
\end{proof}

\begin{theorem}\label{thm:20.8}
If $A$ is a regular UFD then so is $A[X]$.
\end{theorem}

\begin{proof}
Set $B = A[X]$. By Theorem~\hyperref[thm:20.5]{5}, it is enough to prove that $uB \cap vB$
is principal for $u, v \in B$; set $a = uB \cap vB$. Then by Theorem~\hyperref[thm:20.6]{6} and
Theorem~\ref{thm:19.5}, $a$ is projective, so that
\[
  a \otimes_A A = a \otimes_A (B/XB) = a/Xa
\]
is projective as an $A$-module. Suppose that $a = X^rb$ with $b \nsubseteq XB$; then
$a/Xa \cong b/Xb$, so that $b$ is isomorphic to $a$, hence projective, and therefore
locally principal. $B$ is a regular ring, so that the prime divisors of $b$ all
have height 1. Since $XB$ is also a height 1 prime ideal and $b \nsubseteq XB$ we
have $b:XB = b$, hence $b \cap XB = Xb$. Therefore since we can view $b/Xb$
as $b/Xb = b/b \cap XB \subset B/XB = A$, by Theorem~\hyperref[thm:20.7]{7} it is principal, hence
$b = yB + Xb$ for some $y \in b$; then by NAK, $b = yB$, so that $a = X^ryB$.
\end{proof}

\begin{remark}
There are examples where $A$ is a UFD but $A[X]$ is not.
\end{remark}

It is easy to see that a UFD is a Krull ring. For any Krull ring $A$ we
can define the \emph{divisor class group}\index{divisor class group} of $A$, which should be thought of as a
measure of the extent to which $A$ fails to be a UFD. We can give the
definition in simple terms as follows: let $\mathscr{P}$ be the set of height 1 prime
ideals of the Krull ring $A$, and $D(A)$ the free Abelian group on $\mathscr{P}$. That
is, $D(A)$ consists of formal sums $\sum_{p \in \mathscr{P}} n_p \cdot p$ (with $n_p \in \mathbb{Z}$ and all but finitely
many $n_p = 0$), with addition defined by
\[
  \left(\sum n_pp\right) + \left(\sum n_p'p\right) = \sum (n_p + n_p')p.
\]
Let $K$ be the field of fractions of $A$, and $K^*$ the multiplicative group of
non-zero elements of $K$, and for $a \in K^*$ set $\divisor(a) = \sum_{p \in \mathscr{P}} v_p(a) \cdot p$, where $v_p$
is the normalised additive valuation of $K$ corresponding to $p$. Then
$\divisor(ab) = \divisor(a) + \divisor(b)$, so that $\divisor$ is a homomorphism from $K^*$ to $D(A)$. We
write $F(A)$ for the image of $K^*$; this is a subgroup of $D(A)$, so that
we can define $C(A) = D(A)/F(A)$ to be the divisor class group of $A$.

Now, if $A$ is a UFD then each $p \in \mathscr{P}$ is principal, and if $p = aA$ then for
an element of $D(A)$ we have $p = \divisor(a)$, so that $C(A) = 0$. Conversely, if
$C(A) = 0$ then each $p \in \mathscr{P}$ is a principal ideal, and putting this together
with the corollary of Theorem~\ref{thm:12.3}, one sees easily that $A$ is a UFD. Hence
\[
  A \text{ is a UFD} \iff C(A) = 0.
\]

Now let $A$ be any ring, and $M$ a finite projective $A$-module. For each
$P \in \Spec A$, the localisation $M_P$ is a free module over $A_P$, and we write
$n(P)$ for its rank. Then $n$ is a function on $\Spec A$, and is constant on every
connected component (since $n(P) = n(Q)$ if $P \supset Q$). This function $n$ is called
the \emph{rank} of $M$. If the rank is a constant $r$ over the whole of $\Spec A$ then
we say that $M$ is a \emph{projective module of rank $r$}\index{projective module of rank r}. We write $\Pit(A)$ for the
set of isomorphism classes of finite projective $A$-modules of rank 1; $\cl(M)$
denotes the isomorphism class of $M$. If $M$ and $N$ are finite projective rank
1 modules then so is $M \otimes_A N$; this is clear on taking localisations. Thus
we can define a sum in $\Pit(A)$ by setting
\[
  \cl(M) + \cl(N) = \cl(M \otimes N).
\]
We set $M^* = \Hom_A(M, A)$, and define $\varphi\colon M \otimes M^* \to A$ by
\[
  \varphi\left(\sum m_i \otimes f_i\right) = \sum f_i(m_i);
\]
then $\varphi$ is an isomorphism (taking localisations and using the corollary to
Theorem~\ref{thm:7.11} reduces to the case $M = A$, which is clear). Hence $\cl(M^*) =
-\cl(M)$, and $\Pit(A)$ becomes an Abelian group, called the \emph{Picard group}\index{Picard group}
of $A$. If $A$ is local then $\Pit(A) = 0$.

If $A$ is an integral domain with field of fractions $K$, then $M_{(0)} = M \otimes K$,
so that the rank we have just defined coincides with the earlier definition
(after Lemma~\hyperref[lem:20.1]{1}). If $M$ is a finite projective rank 1 module, then since $M$
is torsion-free we have $M \subset M_{(0)} \cong K$, so that $M$ is isomorphic as an
$A$-module to a fractional ideal; for fractional ideals, by Theorem~\ref{thm:11.3},
projective and invertible are equivalent conditions, so that for an integral
domain $A$, we can consider $\Pit(A)$ as a quotient of the group of invertible
fractional ideals under multiplication. A fractional ideal $I$ is isomorphic
to $A$ as an $A$-module precisely when $I$ is principal, so that
\[
  \Pit(A) = (\text{invertible fractional ideals})/(\text{principal fractional ideals}).
\]

Suppose in addition that $A$ is a Krull ring. Then we can view $\Pit(A)$
as a subgroup of $C(A)$. To prove this, for $p \in \mathscr{P}$ and $I$ a fractional ideal,
set
\[
  v_p(I) = \min \{v_p(x) \mid x \in I\};
\]
this is zero for all but finitely many $p \in \mathscr{P}$ (check this!), so that we can set
\[
  \divisor(I) = \sum_{p \in \mathscr{P}} v_p(I) \cdot p \in D(A).
\]
For a principal ideal $I = aA$ we have $\divisor(I) = \divisor(a)$. One sees easily that
$\divisor(II') = \divisor(I) + \divisor(I')$, and that $\divisor(A) = 0$, so that if $I$ is invertible,
$\divisor(I) = -\divisor(I^{-1})$.

For invertible $I$ we have $(I^{-1})^{-1} = I$: indeed, $I \subset (I^{-1})^{-1}$ from the
definition, and $I = I \cdot A \supset I(I^{-1}(I^{-1})^{-1}) = (I^{-1})^{-1}$, so $I = (I^{-1})^{-1}$. If $I$ is invertible and
$\divisor(I) = 0$ then $\divisor(I^{-1}) = 0$, so that $I \subset A$, $I^{-1} \subset A$; hence $A \subset (I^{-1})^{-1} = I$,
and $I = A$. It follows that if $I, I'$ are invertible, $\divisor(I) = \divisor(I')$ implies $I = I'$.
Thus we can view the group of invertible fractional ideals as a subgroup of
$D(A)$, and $\Pit(A)$ as a subgroup of $C(A)$.

If $A$ is a regular ring then as we have seen, $p \in \mathscr{P}$ is a locally free module,
and so is invertible. Clearly from the definition, $\divisor(p) = p$. Hence, in
the case of a regular ring, $D(A)$ is identified with the group of invertible
fractional ideals, and $C(A)$ coincides with $\Pit(A)$.

The notions of $D(A)$ and $\Pit(A)$ originally arise in algebraic geometry.
Let $V$ be an algebraic variety, supposed to be irreducible and normal. We
write $\mathscr{W}$ for the set of irreducible codimension 1 subvarieties of $V$, and
define the group of divisors $D(V)$ of $V$ to be the free Abelian group on $\mathscr{W}$;
a divisor (or Weil divisor) is an element of $D(V)$. Corresponding to a
rational function $f$ on $V$ and an element $W \in \mathscr{W}$, let $v_W(f)$ denote the
order of zero of $f$ along $W$, or minus the order of the pole if $f$ has a pole
along $W$. Write $\divisor(f) = \sum_{W \in \mathscr{W}} v_W(f) \cdot W$ for the divisor of $f$ on $V$ (or just
$\divisor(f)$). For $W \in \mathscr{W}$, the local ring $\mathcal{O}_W$ of $W$ on $V$ is a DVR of the function
field of $V$, and $v_W$ is the corresponding valuation. We say that two divisors
$M, N$ in $D(V)$ are linearly equivalent if their difference $M - N$ is the divisor of a
function, and write $M \sim N$. The quotient group of $D(V)$ by $\sim$, that is the
quotient by the subgroup of divisors of functions, is the divisor class group
of $V$ (up to linear equivalence), and we write $C(V)$ for this. (In addition
to linear equivalence one also considers other equivalence relations with
certain geometric significance (algebraic equivalence, numerical equivalence,
$\ldots$), and divisor class groups, quotients of $D(V)$ by the corresponding subgroups.)

A divisor $M$ on $V$ is said to be a \emph{Cartier divisor} if it is the divisor of a
function in a neighbourhood of every point of $V$. From a Cartier divisor $M$
one constructs a line bundle over $V$, and two Cartier divisors give rise to
isomorphic line bundles if and only if they are linearly equivalent. Cartier
divisors form a subgroup of $D(V)$, and their class group up to linear
equivalence is written $\Pit(V)$; this can also be considered as the group of
isomorphism classes of line bundles over $V$ (with group law defined by
tensor product). If $V$ is smooth then (by Theorem~\hyperref[thm:20.3]{3}) there is no distinction
between Cartier and Weil divisors, and $C(V) = \Pit(V)$.

The reader familiar with algebraic geometry will know that the divisor
class group and Picard group of a Krull ring are an exact translation of
the corresponding notions in algebraic geometry. If $V$ is an affine variety,
with coordinate ring $k[V] = A$ then $C(V) = C(A)$ and $\Pit(V) = \Pit(A)$. In
this case, to say that $A$ is a UFD expresses the fact that every codimension
1 subvariety of $V$ can be defined as the intersection of $V$ with a hypersurface.
If $V \subset \mathbb{P}^n$ is a projective algebraic variety, defined by a prime ideal
$I \subset k[X_0, \ldots, X_n]$, and we set $A = k[X]/I = k[t_0, \ldots, t_n]$ (with $t_i$ the
class of $X_i$) then $A$ is the so-called homogeneous coordinate ring of $V$. If
$A$ is integrally closed we say that $V$ is projectively normal (also
arithmetically normal). This condition is stronger than saying that $V$ is
normal (the local ring of any point of $V$ is normal). If $A$ is a UFD then
every codimension 1 subvariety of $V$ can be given as the intersection in
$\mathbb{P}^n$ of $V$ with a hypersurface. Let $m = (t_0, \ldots, t_n)$ be the homogeneous
maximal ideal of $A$, and write $R = A_m$ for the localisation. The above
statement holds if we just assume that $R$ is a UFD; see Ex.~\ref{ex:20.6-ex}. All the
information about $V$ is contained in the local ring $R$.

Thus $C(A)$, $\Pit(A)$ and the UFD condition are notions with important
geometrical meaning, and methods of algebraic geometry can also be used
in their study. For example, in this way Grothendieck~\cite{G5} was able to
prove the following theorem conjectured by Samuel: let $R$ be a regular
local ring, $P$ a prime ideal generated by an $R$-sequence, and set $A = R/P$;
if $A_p$ is a UFD for every $p \in \Spec A$ with $\htop p \le 3$ then $A$ is a UFD.

We do not have the space to discuss $C(A)$ and $\Pit(A)$ in detail, and we
just mention the following two theorems as examples:
\begin{enumerate}[(1)]
  \item If $A$ is a Krull ring then $C(A) \cong C(A[X])$.
        This generalises the well-known theorem (see Ex.~\ref{ex:20.2-ex}) that if $A$ is a UFD
        then so is $A[X]$.
  \item If $A$ is a regular ring then $C(A) \cong C(A[[X]])$.
        This generalises Theorem~\hyperref[thm:20.8]{8}.
\end{enumerate}
Finally we give an example. Let $k$ be a field of characteristic 0, and set
$A = k[X, Y, Z]/(Z^n - XY) = k[x, y, z]$ for some $n > 1$. Then $A/(z, x) \cong
k[X, Y, Z]/(X, Z) \cong k[Y]$, so that $p = (x, z)$ is a height 1 prime ideal of $A$.
In $D(A)$ we have $np = \divisor(x)$, and it can be proved that $C(A) \cong \mathbb{Z}/d$ (see
\cite{S2}, p.~58). The relation $xy = z^n$ shows that $A$ is not a UFD.

For those wishing to know more about UFDs, consult~\cite{K},~\cite{S2} and
\cite{F}.

% ============ 习题 ============
\section*{Exercises to \S\ref{sec:20}}
\emph{Prove the following propositions.}

\begin{exercise}\label{ex:20.1-ex}
(Gauss' lemma) Let $A$ be a UFD, and $f(X) = a_0 + a_1X + \cdots + a_nX^n \in
A[X]$; say that $f$ is \emph{primitive} if the g.c.d. of the coefficients $a_0, \ldots, a_n$
is 1. Then if $f(X)$ and $g(X)$ are primitive, so is $f(X)g(X)$.
\end{exercise}

\begin{exercise}\label{ex:20.2-ex}
If $A$ is a UFD so is $A[X]$ (use the previous question).
\end{exercise}

\begin{exercise}\label{ex:20.3-ex}
If $A$ is a UFD and $q_1, \ldots, q_r$ are height 1 primary ideals then $q_1 \cap \cdots \cap q_r$ is a
principal ideal.
\end{exercise}

\begin{exercise}\label{ex:20.4-ex}
Let $A$ be a Zariski ring (see \S\ref{sec:8}) and $\hat{A}$ the completion of $A$. Then if $A$ is
a UFD so is $\hat{A}$ (there are counter-examples to the converse).
\end{exercise}

\begin{exercise}\label{ex:20.5-ex}
Let $A$ be an integral domain. We say that $A$ is \emph{locally UFD} if $A_m$ is a UFD
for every maximal ideal $m$. If $A$ is a semilocal integral domain and $A$ is
locally UFD, then $A$ is a UFD.
\end{exercise}

\begin{exercise}\label{ex:20.6-ex}
Let $R = \bigoplus_{n \ge 0} R_n$ be a graded ring, and suppose that $R_0$ is a field. Set
$m = \bigoplus_{n > 0} R_n$. If $I$ is a homogeneous ideal of $R$ such that $IR_m$ is principal
then there is a homogeneous element $f \in I$ such that $I = fR$.
\end{exercise}
